\clearpage
\begin{appendices}
\renewcommand{\thechapter}{\Asbuk{chapter}}
\titleformat{\chapter}[display]{\filcenter}{\MakeUppercase{Приложение} \thechapter}{0pt}{\bfseries}

\chapter{ЛИСТИНГИ ИСХОДНОГО КОДА}

В данном приложении целесообразно разместить фрагменты кода вычислительного
ядра, демонстрирующие реализацию маршрутов \texttt{pt}, \texttt{ph},
\texttt{ps}, \texttt{px} и \texttt{rhot}, а также примеры типизированных
единиц измерения и обработки ошибок.
Дополнительно могут быть приведены выдержки из \texttt{build.rs}, отвечающего за
генерацию таблиц коэффициентов IF97 из исходных CSV-файлов.

\clearpage
\chapter{ДОПОЛНИТЕЛЬНЫЕ ИЛЛЮСТРАТИВНЫЕ МАТЕРИАЛЫ}

В данное приложение рекомендуется вынести дополнительные термодинамические
диаграммы, примеры таблиц результатов расчёта, схему взаимодействия компонентов,
фрагменты \texttt{if97.proto}, а также примеры Docker- и Kubernetes-манифестов.
Размещение таких материалов в приложениях позволяет сохранить компактность
основного текста и при этом зафиксировать инженерные артефакты разработки.

\end{appendices}
